\section{Limitations of Current Approaches}

%SzB: Minden egyes limitation-t alá kell támasztani minimum 1-2 hivatkozással. és azok alapján kell ezeket a gondolatoka megfogalmazni.

Despite significant progress in deep learning-based code translation, current
approaches still face several important limitations that affect their
reliability and applicability in real-world scenarios.

One major issue is hallucination and incorrect logic generation. Models can
produce syntactically valid code that appears plausible but fails to correctly
implement the intended functionality. This highlights the gap between
surface-level correctness and true semantic understanding.

Another limitation is the handling of dependencies and external context. Most
models operate on isolated code snippets and do not fully consider external
libraries, project structure, or runtime dependencies. As a result, generated
code may fail to compile or execute correctly in practical environments.

Context handling also remains a challenge. Although Transformer-based models can
capture long-range dependencies, they are still constrained by fixed context
window sizes. This limitation can lead to incomplete translations or loss of
critical information in larger or more complex codebases.

Furthermore, generalization across datasets and programming languages remains
limited. Models trained on specific datasets often perform well within the same
distribution but struggle when applied to unseen languages, coding styles, or
problem domains.

These limitations are closely connected to challenges in evaluation. Studies
have shown that models achieving high scores on text-based metrics may still
fail in execution-based evaluation, indicating that token-level similarity does
not reliably reflect functional correctness \cite{huang2022exeds}.

These issues demonstrate that current approaches are not yet sufficiently robust
for fully automated code translation. Addressing these challenges requires
improvements in both modeling techniques and evaluation methodologies.